\documentclass[11pt]{article}
\usepackage[T1]{fontenc}
\usepackage{lmodern}
\usepackage[margin=1in]{geometry}
\usepackage{amsmath,amssymb,amsthm,mathtools}
\usepackage{microtype}
\usepackage[hidelinks]{hyperref}

\newtheorem{lemma}{Lemma}[section]
\newtheorem{proposition}[lemma]{Proposition}
\newtheorem{corollary}[lemma]{Corollary}
\theoremstyle{remark}
\newtheorem{remark}[lemma]{Remark}
\newtheorem*{replacement}{Corollary 6.3 (Quantitative stability)}

\newcommand{\ind}{\mathbf{1}}
\newcommand{\E}{\mathbb{E}}
\newcommand{\R}{\mathbb{R}}
\newcommand{\dsq}{\delta_{\square}}
\newcommand{\Tr}{\operatorname{Tr}}
\newcommand{\op}{\mathrm{op}}
\newcommand{\HS}{\mathrm{HS}}
\newcommand{\norm}[1]{\left\lVert #1\right\rVert}
\newcommand{\abs}[1]{\left\lvert #1\right\rvert}

\title{Quantitative stability for the odd-cycle Goodman bound}
\author{}
\date{}

\begin{document}
\maketitle

\begin{abstract}
Assuming the results proved before the stability corollary in
\emph{Goodman-type lower bounds for odd cycles I: The dense regime},
we establish an explicit quantitative replacement for that corollary.
At every density $p_k=1-1/k$, with $k\geq3$, an edge-density error and
an odd-cycle excess of at most $\delta$ force distance
$O_{m,k}(\sqrt{\delta})$ from the balanced complete $k$-partite graphon.
The conclusion holds in relabelled $L^1$ distance, and the exponent
$1/2$ is optimal, even at exact edge density.
The argument consists of a non-degenerate degree estimate, robust
spectral inequalities, and an explicit construction of the parts.
We also record the relaxed qualitative compactness proof.
\end{abstract}

\section{Notation and the quantitative statement}

References to numbered results of the manuscript mean the numbering in
\cite{dense}. In particular, the results below are conditional on its
Theorem 1.1, excursion expansion and polynomial nonnegativity, and
Gamma-moment lemma. We do not assume its qualitative stability corollary
in the quantitative argument.

All graphons are defined on $[0,1]$ with Lebesgue measure. Integrals with
unspecified domains are over the appropriate powers of $[0,1]$.
For a graphon $W$, write
\[
 p=t(K_2,W),\qquad U=1-W,\qquad q=1-p,
 \qquad d_W(x)=\int_0^1W(x,y)\,dy.
\]
The associated integral operator is
\[
 (T_Wf)(x)=\int_0^1W(x,y)f(y)\,dy.
\]
The path $P_j$ has $j$ edges. Throughout, $m\geq3$ is odd, and
\[
 G_m(p)=p^m-p(1-p)^{m-1}.
\]
We use the manuscript's definition
\begin{equation}\label{eq:Phi-definition}
 \Phi_m(W)=t(C_m,W)+t(C_m,U)-t(P_{m-1},U)-G_m(p).
\end{equation}
In particular, the one-colour gap identity is
\begin{equation}\label{eq:gap-decomposition}
 t(C_m,W)-G_m(p)
 =\Phi_m(W)+t(P_{m-1},U)-t(C_m,U).
\end{equation}
For $p\geq2/3$, both terms on the right are nonnegative; see
\cite[Eq.~(6.3)]{dense}.

For a bounded measurable kernel $F$, define
\[
 \norm{F}_{\square}
 =\sup_{S,T\subseteq[0,1]}
   \abs{\int_{S\times T}F(x,y)\,dx\,dy},
\]
where $S,T$ are measurable. Besides the cut distance, define
\[
 \dsq(W,V)=\inf_\phi\norm{W-V^\phi}_{\square},
 \qquad
 \delta_1(W,V)=\inf_\phi\norm{W-V^\phi}_1.
\]
Here $V^\phi(x,y)=V(\phi(x),\phi(y))$, and $\phi$ ranges over
measure-preserving bijections modulo null sets. Directly from these
definitions,
\begin{equation}\label{eq:distance-basics}
 \dsq(W,V)\leq\delta_1(W,V)\leq1,
 \qquad
 \delta_1(1-W,1-V)=\delta_1(W,V).
\end{equation}
Both distances satisfy the triangle inequality: composition of the
relabelings and invariance of the norms under those relabelings reduce
this to the triangle inequality for the respective norms.

For an integer $k\geq3$, set
\[
 q_k=\frac1k,\qquad p_k=1-\frac1k.
\]
Let $W_k$ be the balanced complete $k$-partite graphon, as in the
manuscript, and put $U_k=1-W_k$.

Here are constants used throughout the proof:
\begin{align}
 \alpha_m&=\frac{6m}{7}\left(\frac23\right)^{m-1},\notag\\
 c_m&=\max\left\{2^{4-m},\,
        \alpha_m-(3m+72)2^{-m}\right\},\label{eq:c}\\
 K_k&=64k^4,\label{eq:K}\\
 A_{m,k}&=(22K_k+2)k^{m-2},\label{eq:A}\\
 B_{m,k}&=4K_k+(22K_k+2)k(m-2),\label{eq:B}\\
 L_{m,k}&=1+(m+1)A_{m,k},\label{eq:L}\\
 M_{m,k}&=B_{m,k}\sqrt{\frac{m+1}{c_m}},\label{eq:M}\\
 C_{m,k}&=\sqrt{L_{m,k}}+M_{m,k}.\label{eq:C}
\end{align}
In particular, $c_m>0$.

\begin{replacement}
Fix an odd integer $m\geq3$ and an integer $k\geq3$.
For every $\delta\geq0$, every graphon $W$ satisfying
\begin{equation}\label{eq:hypotheses}
 \abs{t(K_2,W)-p_k}\leq\delta,
 \qquad t(C_m,W)\leq G_m(p_k)+\delta
\end{equation}
satisfies
\begin{equation}\label{eq:main}
 \begin{split}
 \dsq(W,W_k)&\leq\delta_1(W,W_k)\\
 &\leq\min\{1,L_{m,k}\delta+M_{m,k}\sqrt{\delta}\}
 \leq C_{m,k}\sqrt{\delta}.
 \end{split}
\end{equation}
The exponent $1/2$ cannot be replaced by any larger exponent, even
under the additional requirement $t(K_2,W)=p_k$.
\end{replacement}

The proof of the upper bound is completed in Section~\ref{sec:completion},
and optimality is proved in Section~\ref{sec:sharpness}.
All constants are explicit. The exponent and the asymptotic size of the
uniform coefficient $c_m$ are sharp; the remaining numerical constants
and their dependence on $k$ are not optimized.

\section{The relaxed qualitative statement}

\begin{corollary}[Qualitative stability with an edge-density error]
\label{cor:qualitative}
Fix an odd integer $m\geq3$ and an integer $k\geq3$.
For every $\varepsilon>0$, there is $\delta>0$ such that every graphon
satisfying \eqref{eq:hypotheses} has $\dsq(W,W_k)<\varepsilon$.
\end{corollary}

\begin{proof}
Suppose the assertion fails. Then some $\varepsilon_0>0$ has the
following property: for every integer $j\geq1$, there is a graphon
$V_j$ such that
\[
 \abs{t(K_2,V_j)-p_k}\leq j^{-1},\qquad
 t(C_m,V_j)\leq G_m(p_k)+j^{-1},\qquad
 \dsq(V_j,W_k)\geq\varepsilon_0.
\]
The space of graphons modulo zero cut distance is compact, and
homomorphism densities are continuous in cut distance; these are the
facts recalled immediately before \cite[Corollary~6.3]{dense}.
Passing to a subsequence, there is a graphon $V$ with
$\dsq(V_j,V)\to0$. Continuity gives
\[
 t(K_2,V)=p_k,\qquad t(C_m,V)\leq G_m(p_k).
\]
Since $p_k\geq2/3$, \cite[Theorem~1.1]{dense} applies to $V$ and gives
the reverse inequality. Thus $t(C_m,V)=G_m(p_k)$.
By \cite[Proposition~6.2]{dense}, $\dsq(V,W_k)=0$.
The triangle inequality now gives
\[
 \varepsilon_0\leq\dsq(V_j,W_k)
 \leq\dsq(V_j,V)+\dsq(V,W_k)
 =\dsq(V_j,V)\longrightarrow0,
\]
a contradiction.
\end{proof}

\begin{remark}[The endpoint $k=3$]
The preceding proof applies the dense-regime lower bound only to the
limit $V$, whose density is exactly $p_k$. It does not apply that bound
to the approximating graphons $V_j$. Their densities may therefore lie
below $2/3$ when $k=3$.
\end{remark}

\begin{remark}[Finite-graph normalization]
For the usual step graphon $W_G$ associated with a simple graph on $n$
vertices,
\[
 t(K_2,W_G)=\frac{2e(G)}{n^2}
 =\frac{n-1}{n}\frac{e(G)}{\binom n2}.
\]
In particular, when $k\mid n$, the balanced Tur\'an graph has graphon
edge density exactly $p_k$. The relaxed hypothesis is useful because it
allows edge-count errors and the $O(1/n)$ change of normalization; it
should not be motivated by a claim that exact equality is essentially
unavailable for finite graphs.
\end{remark}

\section{A non-degenerate degree estimate}

For $d\geq0$, let $h_d$ be the complete homogeneous symmetric polynomial
of degree $d$, with $h_0=1$, as in \cite[Eq.~(3.1)]{dense}.
For $p=1-q$, the first-excursion polynomial is
\begin{equation}\label{eq:first-poly}
 P_{m,1}(q;\lambda)
 =m\bigl[h_{m-2}(p,-\lambda)+h_{m-2}(q,\lambda)\bigr]
   -h_{m-3}(q,q,\lambda).
\end{equation}

\begin{lemma}[Uniform lower bound for the first excursion]
\label{lem:first-poly}
For $0\leq q\leq1/3$ and $|\lambda|\leq1/2$,
\begin{equation}\label{eq:poly-lower}
 P_{m,1}(q;\lambda)\geq c_m.
\end{equation}
If
\[
 c_m^*=\min_{\substack{0\leq q\leq1/3\\|\lambda|\leq1/2}}
             P_{m,1}(q;\lambda),
\]
then
\begin{equation}\label{eq:c-star-bounds}
 \alpha_m-(3m+72)2^{-m}
 \leq c_m^*\leq\alpha_m+12m2^{-m}.
\end{equation}
Consequently, as $m\to\infty$ through odd integers,
\begin{equation}\label{eq:c-asymptotic}
 c_m^*=\frac{9m}{7}\left(\frac23\right)^m
       \left(1+O\bigl((3/4)^m\bigr)\right),
 \qquad c_m\sim c_m^*.
\end{equation}
The infimum with $0<q\leq1/3$ has the same value, by continuity.
\end{lemma}

\begin{proof}
Set $n=m-2$, which is odd. Expanding \eqref{eq:first-poly} gives
\begin{equation}\label{eq:poly-expansion}
 P_{m,1}(q;\lambda)
 =\sum_{j=0}^{n-1}
 \left[m\bigl((-1)^jp^{n-j}+q^{n-j}\bigr)
       -(n-j)q^{n-1-j}\right]\lambda^j.
\end{equation}
The degree-$n$ terms cancel. For every odd $j$ in this sum the
coefficient is nonpositive, since $p\geq q\geq0$.
It follows that
\begin{equation}\label{eq:reflection}
 P_{m,1}(q;-a)\geq P_{m,1}(q;a)\qquad(a\geq0).
\end{equation}
It is therefore enough to consider $\lambda\geq0$.

\paragraph{A positive bound for every $m$.}
We make explicit the quantitative conclusion contained in the proof of
\cite[Proposition~5.6]{dense}. Define, as in that proof,
\[
 \rho_{n,m}(u)=\frac{m}{n}\bigl(u^n+(1-u)^n\bigr)-u^{n-1},
 \qquad
 a_j=2^{2j+1-n}\binom{n}{2j}
 \quad\left(0\leq j\leq\frac{n-1}{2}\right).
\]
Since $m=n+2$, the centered identity \cite[Eq.~(5.18)]{dense} becomes
\begin{equation}\label{eq:centered}
 n\rho_{n,m}\left(\frac12+v\right)
 =2a_0+\sum_{j=1}^{(n-1)/2}
 a_j\bigl(2(j+1)v^{2j}-jv^{2j-1}\bigr).
\end{equation}
Let $Y\sim\operatorname{Gamma}(1,1)$, with shape and rate as parameters,
and let $x\geq0$. Put
\[
 b=\frac12-q\geq\frac16,\qquad z=x/b,
 \qquad v=q+xY-\frac12=b(zY-1).
\]
The manuscript's Gamma-moment inequality
\cite[Lemma~5.5, with $r=1$]{dense} states that
\[
 3j\E(zY-1)^{2j-1}\leq(j+1)\E(zY-1)^{2j}.
\]
Multiplying by the appropriate powers of $b$ yields
\begin{align*}
 \E\bigl(2(j+1)v^{2j}-jv^{2j-1}\bigr)
 &\geq b^{2j}(j+1)
       \left(2-\frac{1}{3b}\right)\E(zY-1)^{2j}\\
 &\geq0.
\end{align*}
All $a_j$ are positive. Hence \eqref{eq:centered} implies
\begin{equation}\label{eq:strict-gamma}
 n\E\rho_{n,m}(q+xY)\geq2a_0=2^{2-n}.
\end{equation}
This includes $n=1$, when the sum in \eqref{eq:centered} is empty.

For $\lambda>0$, the Gamma representation
\cite[Eq.~(5.15), with $r=1$]{dense} gives
\[
 P_{m,1}(q;\lambda)
 =n\E\rho_{n,m}\left(q+\frac{\lambda Y}{Z}\right),
 \qquad Z\sim\operatorname{Gamma}(m-1,1),
\]
where $Y$ and $Z$ are independent. Conditioning on $Z$ and applying
\eqref{eq:strict-gamma} gives
\begin{equation}\label{eq:baseline}
 P_{m,1}(q;\lambda)\geq2^{2-n}=2^{4-m}.
\end{equation}
These expectations are absolutely integrable: $\rho_{n,m}$ has degree
at most $n-1$, and $Z$ has shape $m-1=n+1$, so the needed inverse
moments of $Z$ are finite. Continuity gives \eqref{eq:baseline} at
$\lambda=0$, and \eqref{eq:reflection} extends it to negative $\lambda$.

\paragraph{An asymptotically sharp bound.}
Assume $0\leq\lambda\leq1/2$. Since $n=m-2$ is odd,
\[
 h_n(p,-\lambda)
 =\frac{p^{m-1}-\lambda^{m-1}}{p+\lambda}.
\]
The function $p\mapsto p^{m-1}/(p+1/2)$ is increasing for $p>0$;
its derivative has numerator
$p^{m-2}((m-2)p+(m-1)/2)>0$.
Therefore
\[
 \frac{mp^{m-1}}{p+\lambda}\geq\alpha_m,
 \qquad
 \frac{m\lambda^{m-1}}{p+\lambda}\leq3m2^{-m}.
\]
Also,
\begin{align*}
 h_{m-3}(q,q,\lambda)
 &=\sum_{j=0}^{m-3}(j+1)q^j\lambda^{m-3-j}\\
 &\leq2^{3-m}\sum_{j=0}^{m-3}(j+1)(2/3)^j
 \leq72\,2^{-m},
\end{align*}
because $\sum_{j\geq0}(j+1)(2/3)^j=9$.
The term $m h_{m-2}(q,\lambda)$ is nonnegative. Thus
\[
 P_{m,1}(q;\lambda)
 \geq\alpha_m-(3m+72)2^{-m}.
\]
Together with \eqref{eq:reflection} and \eqref{eq:baseline}, this proves
\eqref{eq:poly-lower} and the lower bound in \eqref{eq:c-star-bounds}.

At $(q,\lambda)=(1/3,1/2)$, the first positive term
$mp^{m-1}/(p+\lambda)$ equals $\alpha_m$ exactly, whereas
\[
 m h_{m-2}(1/3,1/2)
 \leq m2^{2-m}\sum_{j\geq0}(2/3)^j
 =12m2^{-m}.
\]
Discarding the two nonpositive terms proves the upper bound in
\eqref{eq:c-star-bounds}. Since $\alpha_m=(9m/7)(2/3)^m$,
dividing by $\alpha_m$ proves the asserted asymptotic formula for
$c_m^*$. Moreover,
\[
 \frac{2^{4-m}}{\alpha_m}\longrightarrow0.
\]
The definition of $c_m$ therefore gives $c_m\sim\alpha_m\sim c_m^*$.

\end{proof}

\begin{remark}[The first two lengths]
The polynomial $P_{3,1}$ is identically $2$, and the choice
\eqref{eq:c} gives $c_3=2$.
For $m=5$, the completed-square identity in
\cite[Section~4]{dense} gives
\[
 P_{5,1}(q;\lambda)
 =4\left(\lambda+q-\frac58\right)^2
  +8\left(q-\frac58\right)^2+\frac5{16}.
\]
On the rectangle of Lemma~\ref{lem:first-poly}, its exact minimum is
$143/144$, attained at $(q,\lambda)=(1/3,7/24)$.
The value $c_5=143/144$ may therefore be substituted in all the
stability constants. No assertion about the exact minimizer for
other lengths is needed for the proof.
\end{remark}

\begin{proposition}[Non-degenerate degree estimate]\label{prop:degree}
Every graphon $W$ with edge density $p\geq2/3$ satisfies
\begin{equation}\label{eq:degree-coercivity}
 \Phi_m(W)\geq c_m\int_0^1(d_W(x)-p)^2\,dx.
\end{equation}
\end{proposition}

\begin{proof}
First suppose $W$ is a step graphon as in \cite[Section~2]{dense}.
The excursion expansion \cite[Eq.~(3.10)]{dense} reads
\[
 \Phi_m(W)=\sum_{r=1}^{(m-1)/2}
 \int P_{m,r}(q;\lambda_1,\ldots,\lambda_r)
       \,d\mu(\lambda_1)\cdots d\mu(\lambda_r).
\]
Here $\mu$ is nonnegative, is supported on $[-1/2,1/2]$, and satisfies
\[
 \mu(\R)=\norm{g}_2^2
 =\int_0^1(d_W(x)-p)^2\,dx;
\]
see \cite[Eq.~(2.3) and the proof of Proposition~6.1]{dense}.
All excursion integrands are nonnegative by
\cite[Section~5]{dense}. Retaining only $r=1$ and applying
Lemma~\ref{lem:first-poly} proves \eqref{eq:degree-coercivity}.

For a general graphon, let $W_s$ be its same-density dyadic step
approximations from \cite[Lemma~2.1 and Appendix~A]{dense}.
Then $\norm{W_s-W}_1\to0$, and every homomorphism density in
\eqref{eq:Phi-definition} converges. Also,
\[
 \int(d_{W_s}-p)^2=t(P_2,W_s)-p^2
 \longrightarrow t(P_2,W)-p^2=\int(d_W-p)^2.
\]
Passing to the limit proves the assertion.
\end{proof}

\section{Robust spectral estimates}

For any graphon $U$, define
\begin{align}
 \eta(U)&=\int U(x,y)(1-U(x,y))\,dx\,dy,\label{eq:eta-def}\\
 \tau(U)&=t(P_2,U)-t(C_3,U)\notag\\
 &=\int U(x,y)U(y,z)(1-U(x,z))\,dx\,dy\,dz.\label{eq:tau-def}
\end{align}
Both are nonnegative.

\begin{lemma}[Even-path inequality]\label{lem:path}
If $U$ has edge density $q$ and $\ell\geq2$ is even, then
$t(P_\ell,U)\geq q^\ell$.
\end{lemma}

\begin{proof}
Let $\nu$ be the spectral measure of the self-adjoint operator $T_U$
at the unit vector $\ind$. Thus $\nu$ is a probability measure and
\[
 \int\lambda\,d\nu(\lambda)=\langle\ind,T_U\ind\rangle=q,
 \qquad
 \int\lambda^\ell\,d\nu(\lambda)
 =\langle\ind,T_U^\ell\ind\rangle=t(P_\ell,U).
\]
The function $x\mapsto x^\ell$ is convex on $\R$.
Jensen's inequality proves the claim.
\end{proof}

\begin{lemma}[Robust moment comparison]\label{lem:spectral}
Let $U$ be a graphon of edge density $q>0$, and put
\[
 \sigma^2=\int_0^1(d_U-q)^2.
\]
Suppose $\Gamma\geq0$ and
\begin{equation}\label{eq:moment-hypothesis}
 t(C_m,U)\geq q^{m-1}-\Gamma.
\end{equation}
Then
\begin{align}
 \norm{T_U}_{\op}
 &\leq q+\frac{\sigma}{q},\label{eq:operator-bound}\\
 \eta(U)
 &\leq q^{2-m}\Gamma+\frac{m-2}{q}\sigma,\label{eq:eta-bound}\\
 \tau(U)
 &\leq\sigma^2+2q^{3-m}\Gamma+2(m-2)\sigma.\label{eq:tau-bound}
\end{align}
Moreover,
\begin{equation}\label{eq:HS-bound}
 \norm{T_U^2-qT_U}_{\HS}^2
 \leq8q^{3-m}\Gamma+8(m-2)\sigma+\frac{2\sigma^2}{q}.
\end{equation}
\end{lemma}

\begin{proof}
Write $a=m-2$, a positive odd integer, and
$\varrho=\norm{T_U}_{\op}$. Since $T_U$ is Hilbert--Schmidt and
self-adjoint, it is compact and has real spectrum. Also,
\[
 q=\langle\ind,T_U\ind\rangle\leq\varrho
 \leq\norm{U}_2\leq1.
\]
There is a nonnegative unit eigenfunction $f$ with eigenvalue
$\varrho$. To verify this without a positivity theorem, take a real
unit eigenfunction $h$ for an eigenvalue of modulus $\varrho$.
Nonnegativity of $U$ gives
\[
 \varrho=\abs{\langle h,T_Uh\rangle}
 \leq\langle |h|,T_U|h|\rangle\leq\varrho.
\]
Thus $f=|h|$ attains the upper operator-norm bound in its quadratic
form. The identity
\[
 \norm{T_Uf-\varrho f}_2^2
 =\norm{T_Uf}_2^2-2\varrho\langle f,T_Uf\rangle+\varrho^2
 \leq0
\]
shows that $T_Uf=\varrho f$.

Since $U\leq1$ and $\norm{f}_2=1$,
\[
 \norm{f}_\infty
 \leq\varrho^{-1}\norm{f}_1\leq\varrho^{-1},
 \qquad \norm{f^2}_2\leq\norm{f}_\infty\norm{f}_2
 \leq\varrho^{-1}.
\]
Using symmetry and $2f(x)f(y)\leq f(x)^2+f(y)^2$ gives
\begin{align*}
 \varrho
 &=\int U(x,y)f(x)f(y)\,dx\,dy\\
 &\leq\int d_U(x)f(x)^2\,dx
 \leq q+\sigma\norm{f^2}_2
 \leq q+\frac{\sigma}{\varrho}
 \leq q+\frac{\sigma}{q}.
\end{align*}
This proves \eqref{eq:operator-bound}.

Let $(\lambda_i)$ list the nonzero eigenvalues with multiplicity.
The Hilbert--Schmidt and cycle trace identities give
\[
 \sum_i\lambda_i^2=\int U^2\leq q,
 \qquad t(C_m,U)=\sum_i\lambda_i^m.
\]
All series used below converge absolutely, because
$|\lambda_i|\leq1$ and $\sum_i\lambda_i^2<\infty$.
Since $a$ is odd and $\lambda_i\leq\varrho$, each summand in
\begin{equation}\label{eq:F}
 F:=\varrho^a q-t(C_m,U)
 =\varrho^a\eta(U)
  +\sum_i\lambda_i^2(\varrho^a-\lambda_i^a)
\end{equation}
is nonnegative. The hypothesis implies
\begin{equation}\label{eq:F-upper}
 0\leq F\leq\Gamma+q(\varrho^a-q^a).
\end{equation}
For $0\leq x\leq1$,
$1-x^a\leq a(1-x)$, by factoring $1-x^a$.
Consequently,
\begin{align*}
 \eta(U)
 &\leq\varrho^{-a}\Gamma
       +q\bigl(1-(q/\varrho)^a\bigr)\\
 &\leq q^{-a}\Gamma+a(\varrho-q)
 \leq q^{-a}\Gamma+\frac{a}{q}\sigma.
\end{align*}
This is \eqref{eq:eta-bound}.

For $-1\leq x\leq1$, oddness of $a$ gives
\begin{equation}\label{eq:power-comparison}
 1-x\leq2(1-x^a).
\end{equation}
Indeed, for $x\geq0$ one has $x^a\leq x$, and for $x\leq0$
the left side is at most $2$ whereas $1-x^a\geq1$.
Using \eqref{eq:power-comparison} in the spectral sum gives
\begin{align*}
 \varrho q-t(C_3,U)
 &=\varrho\eta(U)+\sum_i\lambda_i^2(\varrho-\lambda_i)\\
 &\leq2\varrho^{1-a}F.
\end{align*}
Since $t(P_2,U)=q^2+\sigma^2$ and $\varrho\geq q$,
\begin{align*}
 \tau(U)
 &=\sigma^2+q^2-\varrho q+\varrho q-t(C_3,U)\\
 &\leq\sigma^2+2\varrho^{1-a}\Gamma
       +2\varrho q\bigl(1-(q/\varrho)^a\bigr)\\
 &\leq\sigma^2+2q^{1-a}\Gamma+2a q(\varrho-q)\\
 &\leq\sigma^2+2q^{1-a}\Gamma+2a\sigma.
\end{align*}
Here $\varrho^{1-a}\leq q^{1-a}$ because $a\geq1$.
This proves \eqref{eq:tau-bound}.

Finally, put
\[
 S=\sum_i\lambda_i^2(\varrho-\lambda_i)\geq0.
\]
The same argument, omitting the nonnegative term
$\varrho\eta(U)$, gives
\[
 S\leq2q^{1-a}\Gamma+2a\sigma.
\]
Since $0\leq\varrho-\lambda_i\leq2\varrho$, we obtain
\begin{align*}
 \norm{T_U^2-qT_U}_{\HS}^2
 &=\sum_i\lambda_i^2(\lambda_i-q)^2\\
 &\leq2\sum_i\lambda_i^2(\lambda_i-\varrho)^2
       +2(\varrho-q)^2\sum_i\lambda_i^2\\
 &\leq4\varrho S+2q(\varrho-q)^2\\
 &\leq8q^{1-a}\Gamma+8a\sigma+2\sigma^2/q,
\end{align*}
using $\varrho\leq1$ and \eqref{eq:operator-bound}.
This is \eqref{eq:HS-bound}.
\end{proof}

\begin{remark}
The deficit in \eqref{eq:F} contains the signed odd power
$\lambda_i^{m-2}$, not $|\lambda_i|^{m-2}$.
Negative eigenvalues therefore contribute positively to that deficit.
Although \eqref{eq:HS-bound} is a robust version of the manuscript's
projection identity, we will not take its square root to estimate the
structural distance. Instead, we use \eqref{eq:eta-bound} and
\eqref{eq:tau-bound} directly in the following partition construction.
\end{remark}

\section{Recovering a balanced clique partition}\label{sec:rounding}

\begin{lemma}[Approximation by a balanced union of cliques]
\label{lem:rounding}
Let $k\geq2$, let $q=1/k$, and let $H$ be a $\{0,1\}$-valued graphon.
Set
\[
 v=\int_0^1|d_H-q|,\qquad
 \tau=\int H(x,y)H(y,z)(1-H(x,z))\,dx\,dy\,dz.
\]
Then
\begin{equation}\label{eq:rounding}
 \delta_1(H,U_k)
 \leq16k^2\bigl(v+12k^2\tau\bigr)
 \leq64k^4(v+3\tau).
\end{equation}
\end{lemma}

\begin{proof}
Write $b=v+12k^2\tau$. If $b\geq1/(16k^2)$, the first inequality
in \eqref{eq:rounding} follows from $\delta_1\leq1$.
Henceforth assume
\begin{equation}\label{eq:small-b}
 b<\frac{1}{16k^2}.
\end{equation}
We may replace $H$ on a null set by a symmetric Borel
$\{0,1\}$-valued representative. All sets used below may then be
taken Borel, and all integral identities are unaffected.

\paragraph{The cost of selecting a neighbourhood.}
For a remaining measurable set $R\subseteq[0,1]$ and $x\in R$, define
\[
 N_R(x)=\{y\in R:H(x,y)=1\},\qquad d_R(x)=|N_R(x)|.
\]
For a measurable set $A$, $|A|$ denotes its Lebesgue measure.
Define
\begin{equation}\label{eq:local-cost}
 e_R(x)=\int_{N_R(x)^2}(1-H)
       +2\int_{N_R(x)\times(R\setminus N_R(x))}H.
\end{equation}
This is exactly the $L^1$ cost of making $N_R(x)$ a clique and
removing its edges to the rest of $R$.
All costs are computed against the original graphon $H$; it is not
modified during the selection process.

Let $\tau_R$ be the integral defining $\tau$, restricted to $R^3$.
Substituting the indicator $H(x,\cdot)$ for membership in $N_R(x)$,
Fubini's theorem and symmetry give
\begin{align}
 \int_R e_R(x)\,dx
 &=\int_{R^3}H(x,y)H(x,z)(1-H(y,z))\,dx\,dy\,dz\notag\\
 &\quad+2\int_{R^3}H(x,y)(1-H(x,z))H(y,z)\,dx\,dy\,dz\notag\\
 &=3\tau_R\leq3\tau.\label{eq:average-cost}
\end{align}

\paragraph{Extraction with a uniform cost bound.}
Suppose disjoint sets $A_1,\ldots,A_t$ have been selected, and let
$R=[0,1]\setminus\bigcup_{i=1}^t A_i$ be the remainder.
Let $E$ be the sum of their costs \eqref{eq:local-cost} at the times
when they were selected. Every original edge from $R$ to a selected
set was counted with its symmetric counterpart at that selection.
All other contributions to $E$ are nonnegative. Thus
\[
 \int_{R\times R^c}H\leq E/2.
\]
Since $d_H(x)-d_R(x)=\int_{R^c}H(x,y)\,dy$ on $R$, we have
\begin{equation}\label{eq:remaining-degree}
 \int_R|d_R-q|
 \leq v+\int_{R\times R^c}H
 \leq v+E/2.
\end{equation}

We select each set to have measure at least $q/2$ and cost at most
$12k\tau$. Suppose inductively that $t$ such choices have been made.
Their disjointness gives $t\leq2k$, so
\begin{equation}\label{eq:total-cost-bound}
 E\leq12kt\tau\leq24k^2\tau,
 \qquad v+E/2\leq b.
\end{equation}
If $|R|\geq q/2$, define
\[
 X_R=\{x\in R:d_R(x)\geq q/2\}.
\]
On $R\setminus X_R$, one has $|d_R-q|\geq q/2$, whence
\[
 |R\setminus X_R|\leq2k(v+E/2)\leq2kb<q/8.
\]
In particular, $|X_R|\geq q/4$.
By \eqref{eq:average-cost}, there is $x\in X_R$ such that
\[
 e_R(x)\leq\frac{3\tau}{|X_R|}\leq12k\tau.
\]
Choose $A_{t+1}=N_R(x)$ and replace $R$ by $R\setminus A_{t+1}$.
This preserves the asserted lower bound on the selected measures and
the upper bound on each cost. The induction starts with no selected
sets and $E=0$.

Each choice removes measure at least $q/2$, so the process terminates
after at most $2k$ choices. Write the final sets as
$A_1,\ldots,A_N$, put $a_i=|A_i|$, and denote the final remainder by
$R$, of measure $r$. Then
\begin{equation}\label{eq:extraction-output}
 a_i\geq q/2,\qquad N\leq2k,\qquad r<q/2,
 \qquad E\leq24k^2\tau.
\end{equation}

\paragraph{The remainder and the number of parts.}
For $x\in R$, $d_R(x)\leq r<q/2$. Thus
\eqref{eq:remaining-degree} gives
\[
 \frac{qr}{2}\leq\int_R|d_R-q|\leq b,
 \qquad r\leq2kb.
\]
It also follows that $r^2\leq b$, because $r<q/2$.
Define
\[
 V_0=\sum_{i=1}^N\ind_{A_i\times A_i},
\]
with value zero on pairs involving the remainder unless both endpoints
belong to one of the selected sets. The costs counted during extraction
account exactly for all discrepancies outside $R^2$. Therefore
\begin{equation}\label{eq:partial-cliques}
 \norm{H-V_0}_1=E+\int_{R^2}H\leq E+r^2\leq3b.
\end{equation}
The degree function of $V_0$ equals $a_i$ on $A_i$ and zero on $R$.
Since $\norm{d_H-d_{V_0}}_1\leq\norm{H-V_0}_1$, the more precise
bound $v+E+r^2\leq3b$ yields
\begin{equation}\label{eq:part-degrees}
 \sum_{i=1}^N a_i|a_i-q|+qr
 =\int|d_{V_0}-q|
 \leq v+E+r^2\leq3b.
\end{equation}
Here $E\leq24k^2\tau=2(b-v)$ was used in the last inequality.
As $a_i\geq1/(2k)$, \eqref{eq:part-degrees} implies
\[
 \sum_{i=1}^N|a_i-q|\leq6kb.
\]
Since $\sum_i a_i=1-r$, we conclude
\[
 |1-Nq|\leq r+\sum_i|a_i-q|\leq8kb<q/2.
\]
The quantity $1-Nq=(k-N)/k$ is an integer multiple of $q$.
Consequently,
\begin{equation}\label{eq:number-of-parts}
 N=k.
\end{equation}

\paragraph{Balancing the parts.}
For each $i$, retain a measurable subset $A_i'\subseteq A_i$ with
$|A_i'|=\min(a_i,q)$. This is possible by atomlessness.
The pool of unretained points has measure
\[
 s=r+\sum_{i=1}^k(a_i-q)_+
   =\sum_{i=1}^k(q-a_i)_+
   \leq r+\sum_{i=1}^k|a_i-q|\leq8kb,
\]
where $x_+=\max\{x,0\}$.
The equality follows from $\sum_i a_i=1-r$ and $kq=1$.
Partition this pool into sets of measures $(q-a_i)_+$ and adjoin the
corresponding set to $A_i'$. This gives a measurable partition
$B_1,\ldots,B_k$ of $[0,1]$ with $|B_i|=q$.

Let $V_B=\sum_i\ind_{B_i\times B_i}$.
On pairs whose two endpoints were retained, $V_B$ and $V_0$ agree.
Hence their difference is supported on pairs with at least one
endpoint in the pool, so
\[
 \norm{V_B-V_0}_1\leq2s\leq16kb.
\]
The graphon $V_B$ is a relabelling of $U_k$.
For completeness, on $B_i$ use the map
\[
 \phi(x)=(i-1)q+\int_0^x\ind_{B_i}(t)\,dt.
\]
The cumulative integral is continuous and sends Lebesgue measure
restricted to $B_i$ to Lebesgue measure on $[0,q]$:
for $0\leq u\leq q$, the set of points of $B_i$ at which that
integral is at most $u$ has measure $u$.
The map is strictly increasing on the density-one points of $B_i$,
because between two such distinct points $B_i$ has positive measure.
After discarding null sets it therefore has a measurable inverse,
given by the corresponding generalized inverse. Combining the maps
on the $k$ parts gives a measure-preserving bijection modulo null
sets, and $U_k^\phi=V_B$ almost everywhere.

Combining this with \eqref{eq:partial-cliques}, we obtain
\[
 \delta_1(H,U_k)\leq\norm{H-V_B}_1
 \leq(16k+3)b\leq16k^2b,
\]
where the last inequality holds for $k\geq2$.
Finally,
\[
 16k^2(v+12k^2\tau)\leq64k^4(v+3\tau),
\]
which proves \eqref{eq:rounding}.
\end{proof}

\section{Proof of quantitative stability}\label{sec:completion}

We first obtain an estimate at exact density, with separate dependence
on the cycle gap and the degree deviation.

\begin{proposition}[Exact-density estimate]\label{prop:exact}
Suppose $t(K_2,W)=p_k$, and define
\[
 \Gamma=t(C_m,W)-G_m(p_k),\qquad
 \sigma^2=\int_0^1(d_W-p_k)^2.
\]
Then $\Gamma\geq0$ and
\begin{equation}\label{eq:exact-estimate}
 \delta_1(W,W_k)
 \leq A_{m,k}\Gamma+B_{m,k}\sigma
 \leq A_{m,k}\Gamma+B_{m,k}\sqrt{\Gamma/c_m}.
\end{equation}
\end{proposition}

\begin{proof}
Put $p=p_k$, $q=1/k$, and $U=1-W$.
Then $d_U-q=-(d_W-p)$, so $\sigma^2=\int(d_U-q)^2$.
Theorem 1.1 of the manuscript gives $\Gamma\geq0$.
By \eqref{eq:gap-decomposition} and Proposition~\ref{prop:degree},
\begin{equation}\label{eq:two-small-gaps}
 \sigma^2\leq\Gamma/c_m,
 \qquad 0\leq t(P_{m-1},U)-t(C_m,U)\leq\Gamma.
\end{equation}
Lemma~\ref{lem:path}, with the even length $m-1$, now gives
\[
 t(C_m,U)\geq q^{m-1}-\Gamma.
\]
Thus Lemma~\ref{lem:spectral} applies to $U$.

Define a $\{0,1\}$-valued graphon by thresholding:
\[
 H=\ind_{\{U\geq1/2\}},\qquad e=\norm{H-U}_1.
\]
For $0\leq u\leq1$,
$\min(u,1-u)\leq2u(1-u)$. Hence \eqref{eq:eta-bound} implies
\begin{equation}\label{eq:threshold-error}
 e\leq2\eta(U)
 \leq2k^{m-2}\Gamma+2k(m-2)\sigma.
\end{equation}
Also,
\begin{equation}\label{eq:threshold-degrees}
 v_H:=\int|d_H-q|\leq\sigma+e,
\end{equation}
because $\norm{d_H-d_U}_1\leq e$ and
$\norm{d_U-q}_1\leq\sigma$.
Telescoping the three factors in \eqref{eq:tau-def} gives
\begin{equation}\label{eq:threshold-triples}
 |\tau(H)-\tau(U)|\leq3e.
\end{equation}
Indeed, each of the three differences is bounded in absolute value by
$|H-U|$ on its corresponding pair of coordinates; integration over the
third coordinate contributes a factor of one.

Using Lemma~\ref{lem:rounding}, the triangle inequality, and
\eqref{eq:distance-basics}, we obtain
\begin{align*}
 \delta_1(W,W_k)=\delta_1(U,U_k)
 &\leq e+K_k\bigl(v_H+3\tau(H)\bigr)\\
 &\leq e+K_k\bigl(\sigma+10e+3\tau(U)\bigr).
\end{align*}
Substitute \eqref{eq:threshold-error} and \eqref{eq:tau-bound}:
\begin{align*}
 \delta_1(W,W_k)
 &\leq\bigl[2(1+10K_k)k^{m-2}+6K_k k^{m-3}\bigr]\Gamma\\
 &\quad+\bigl[2(1+10K_k)k(m-2)+6K_k(m-2)+K_k\bigr]\sigma
       +3K_k\sigma^2.
\end{align*}
Since $|d_U-q|\leq1$, we have $0\leq\sigma\leq1$, so
$\sigma^2\leq\sigma$.
Since $k\geq3$, the coefficient of $\Gamma$ is at most
\[
 (22K_k+2)k^{m-2}=A_{m,k},
\]
and the resulting coefficient of $\sigma$ is at most
\[
 (22K_k+2)k(m-2)+4K_k=B_{m,k}.
\]
This proves the first inequality in \eqref{eq:exact-estimate}.
The second follows from \eqref{eq:two-small-gaps}.
\end{proof}

\begin{proof}[Proof of the quantitative corollary]
Let $W$ satisfy \eqref{eq:hypotheses}, and put $p_0=t(K_2,W)$.
Define
\begin{equation}\label{eq:density-correction}
 \widehat W=
 \begin{cases}
 \dfrac{p_k}{p_0}W,&p_0\geq p_k,\\[6pt]
 W+\dfrac{p_k-p_0}{1-p_0}(1-W),&p_0<p_k.
 \end{cases}
\end{equation}
The denominators in the relevant cases are positive, and the
coefficients lie in $[0,1]$, so $\widehat W$ is a graphon.
Direct integration shows that
\begin{equation}\label{eq:density-correction-error}
 t(K_2,\widehat W)=p_k,
 \qquad\norm{\widehat W-W}_1=|p_0-p_k|\leq\delta.
\end{equation}

For graphons $F_1,F_2$ and a finite simple graph $F$, telescoping the
edge factors in the defining integral gives
\[
 |t(F,F_1)-t(F,F_2)|\leq |E(F)|\norm{F_1-F_2}_1.
\]
Applying this to $C_m$, and then applying the dense-regime lower
bound to $\widehat W$, yields
\[
 0\leq\widehat\Gamma
 :=t(C_m,\widehat W)-G_m(p_k)
 \leq(m+1)\delta.
\]
Only $\widehat W$, whose density is exactly $p_k\geq2/3$, is subjected
to the dense-regime theorem and the degree estimate. In particular,
this argument remains valid when $k=3$ and $p_0<2/3$.

Proposition~\ref{prop:exact} and
\eqref{eq:density-correction-error} give
\begin{align*}
 \delta_1(W,W_k)
 &\leq\delta+A_{m,k}\widehat\Gamma
          +B_{m,k}\sqrt{\widehat\Gamma/c_m}\\
 &\leq L_{m,k}\delta+M_{m,k}\sqrt{\delta}.
\end{align*}
Also $\delta_1(W,W_k)\leq1$.
For all $x\geq0$, $\min(1,x)\leq\sqrt{x}$, and for $x,y\geq0$,
$\min(1,x+y)\leq\min(1,x)+y$.
Therefore
\[
 \min\{1,L_{m,k}\delta+M_{m,k}\sqrt{\delta}\}
 \leq\bigl(\sqrt{L_{m,k}}+M_{m,k}\bigr)\sqrt{\delta}.
\]
Together with $\dsq\leq\delta_1$, this proves \eqref{eq:main}.
Optimality of the exponent follows from the next section.
\end{proof}

\section{Optimality of the exponent}\label{sec:sharpness}

Fix $m,k$ and write $q=1/k$, $p=1-q$.
For $0<\varepsilon<q$, let $W_\varepsilon$ be the complete
$k$-partite graphon with part measures
\[
 q+\varepsilon,\quad q-\varepsilon,\quad q,\ldots,q.
\]
Its edge density is
\begin{equation}\label{eq:unbalanced-edge}
 t(K_2,W_\varepsilon)
 =1-\bigl((q+\varepsilon)^2+(q-\varepsilon)^2+(k-2)q^2\bigr)
 =p-2\varepsilon^2.
\end{equation}

\begin{lemma}[Cycle density under a change of part sizes]
\label{lem:expansion}
Define
\begin{equation}\label{eq:beta}
 \beta_{m,k}
 =2mp\bigl(p^{m-1}-q^{m-1}\bigr)
  +m(m-1)(1-2q)q^{m-2}>0.
\end{equation}
For $0<\varepsilon\leq q/2$,
\begin{equation}\label{eq:unbalanced-cycle}
 t(C_m,W_\varepsilon)
 =G_m(p)-\beta_{m,k}\varepsilon^2+R_{m,k}(\varepsilon),
 \qquad |R_{m,k}(\varepsilon)|\leq3^m\varepsilon^4.
\end{equation}
\end{lemma}

\begin{proof}
Let $J$ be the $k\times k$ all-one matrix, let $I$ be the identity,
and put
\[
 D=\operatorname{diag}(1,-1,0,\ldots,0),\qquad
 B_0=q(J-I),\qquad R=(J-I)D.
\]
For part measures $s_1,\ldots,s_k$, summing over the successive part
indices of a cycle gives
\[
 t(C_m,W_{(s_1,\ldots,s_k)})
 =\Tr\!\left(\bigl((J-I)\operatorname{diag}(s_1,\ldots,s_k)\bigr)^m\right).
\]
Consequently,
\[
 t(C_m,W_\varepsilon)=\Tr\!\left((B_0+\varepsilon R)^m\right).
\]
This polynomial is even in $\varepsilon$, because interchanging the
first two part indices sends $\varepsilon$ to $-\varepsilon$.
Its constant term is $G_m(p)$ by \cite[Eq.~(1.1)]{dense}.

The coefficient of $\varepsilon^2$ is
\begin{equation}\label{eq:quadratic-trace}
 \frac m2\sum_{j=0}^{m-2}
 \Tr\bigl(RB_0^jRB_0^{m-2-j}\bigr).
\end{equation}
To see the factor $m/2$, count ordered choices of the two factors
from which $\varepsilon R$ is selected. There are $m$ choices for
the first position, and $j\in\{0,\ldots,m-2\}$ records the number
of intervening $B_0$ factors. Cyclicity of trace makes the value
independent of the first position. Each unordered pair of positions
is counted twice.

Set $P_0=J/k$ and $P_1=I-P_0$. These are complementary orthogonal
projections, and
\[
 B_0=pP_0-qP_1,\qquad P_0DP_0=0,
 \qquad\Tr(P_0D^2)=2q,\qquad\Tr(D^2)=2.
\]
Using $J-I=(k-1)P_0-P_1$, one obtains
\begin{align*}
 P_0RP_0&=0,\\
 \Tr(P_0RP_1R)
 &=-(k-1)\Tr(P_0DP_1D)=-2p,\\
 \Tr(P_1RP_1R)
 &=\Tr(P_1DP_1D)=2-4q=2(1-2q).
\end{align*}
Write $a=m-2$. Expanding the trace in the two projections now gives
\begin{align*}
 \Tr(RB_0^jRB_0^{a-j})
 &=-2p\bigl(p^j(-q)^{a-j}+(-q)^jp^{a-j}\bigr)\\
 &\quad+2(1-2q)(-q)^a.
\end{align*}
Since $a$ is odd and $p+q=1$,
\[
 \sum_{j=0}^a p^j(-q)^{a-j}
 =\frac{p^{a+1}-(-q)^{a+1}}{p+q}
 =p^{m-1}-q^{m-1}.
\]
Substituting in \eqref{eq:quadratic-trace} gives precisely
$-\beta_{m,k}$. Positivity of $\beta_{m,k}$ follows from
$p>q$ and $1-2q>0$.

It remains to bound the higher coefficients.
Express the cycle density as the sum, over proper assignments of
$k$ part indices to the cycle vertices, of the product of their
part measures. For a prescribed set of $j$ vertices contributing
an $\varepsilon$ factor, each such vertex has at most two part
indices with nonzero contribution, of absolute value one.
Each remaining vertex has $k$ choices and contributes a factor $q$.
Dropping the properness restrictions therefore bounds the absolute
coefficient contributed by that vertex set by
\[
 2^j k^{m-j}q^{m-j}=2^j.
\]
The absolute value of the total degree-$j$ coefficient is consequently
at most $\binom mj2^j$. Odd coefficients vanish. Since
$\varepsilon\leq q/2\leq1$, all terms of degree at least four have
absolute sum at most
\[
 \varepsilon^4\sum_{j=4}^m\binom mj2^j
 \leq3^m\varepsilon^4.
\]
This proves \eqref{eq:unbalanced-cycle}.
\end{proof}

\begin{proposition}[Sharpness, including exact edge density]
\label{prop:sharpness}
No constants $C=C(m,k)<\infty$ and $c>1/2$ can replace the
square-root bound in \eqref{eq:main}, even if the bound is required
only for sufficiently small $\delta>0$.
The same assertion holds under the exact-density constraint
$t(K_2,W)=p_k$.
\end{proposition}

\begin{proof}
For every relabelling $W_k^\phi$, the degree function is the constant
$p$. If $S$ is the part of $W_\varepsilon$ of measure
$q+\varepsilon$, then $d_{W_\varepsilon}=p-\varepsilon$ on $S$.
Testing the cut norm on $S\times[0,1]$ gives, for every $\phi$,
\[
 \norm{W_\varepsilon-W_k^\phi}_{\square}
 \geq\abs{\int_S(d_{W_\varepsilon}-p)}
 =(q+\varepsilon)\varepsilon.
\]
Taking the infimum yields
\begin{equation}\label{eq:distance-lower}
 \dsq(W_\varepsilon,W_k)\geq(q+\varepsilon)\varepsilon.
\end{equation}
There is also an explicit upper bound
\begin{equation}\label{eq:distance-upper}
 \delta_1(W_\varepsilon,W_k)
 \leq4q\varepsilon-2\varepsilon^2.
\end{equation}
Indeed, start with two balanced parts of measure $q$ and move a set
of measure $\varepsilon$ from the second to the first.
Only its pairs with the original first part and the remaining
second part change, at total $L^1$ cost
$2q\varepsilon+2\varepsilon(q-\varepsilon)$.
Thus both distances are of order $\varepsilon$.

Set
\[
 \varepsilon_0=
 \min\left\{\frac q2,
       \sqrt{\frac{\beta_{m,k}}{2\cdot3^m}}\right\}>0.
\]
For $0<\varepsilon\leq\varepsilon_0$,
Lemma~\ref{lem:expansion} gives
\[
 t(C_m,W_\varepsilon)
 \leq G_m(p)-\tfrac12\beta_{m,k}\varepsilon^2
 \leq G_m(p).
\]
Together with \eqref{eq:unbalanced-edge}, this shows that the relaxed
hypotheses hold with $\delta=2\varepsilon^2$.
But \eqref{eq:distance-lower} gives distance at least $q\varepsilon$.
For $c>1/2$, the ratio
\[
 \frac{q\varepsilon}{(2\varepsilon^2)^c}
 =q2^{-c}\varepsilon^{1-2c}
\]
tends to infinity as $\varepsilon\downarrow0$.
This proves the first assertion.

For exact density, define
\[
 \widetilde W_\varepsilon
 =W_\varepsilon+
   \frac{2\varepsilon^2}{q+2\varepsilon^2}(1-W_\varepsilon).
\]
Because $\int(1-W_\varepsilon)=q+2\varepsilon^2$, we have
\[
 t(K_2,\widetilde W_\varepsilon)=p,
 \qquad
 \norm{\widetilde W_\varepsilon-W_\varepsilon}_1
 =2\varepsilon^2.
\]
The cycle-density telescoping bound and the preceding estimate give
\[
 t(C_m,\widetilde W_\varepsilon)
 \leq G_m(p)+2m\varepsilon^2.
\]
The triangle inequality, \eqref{eq:distance-lower}, and
$\norm{\cdot}_{\square}\leq\norm{\cdot}_1$ give
\begin{align*}
 \dsq(\widetilde W_\varepsilon,W_k)
 &\geq(q+\varepsilon)\varepsilon-2\varepsilon^2\\
 &=q\varepsilon-\varepsilon^2
 \geq(q/2)\varepsilon.
\end{align*}
Now use $\delta=2m\varepsilon^2$ and let $\varepsilon\downarrow0$.
This again rules out every $c>1/2$.
\end{proof}

\section{Dependence on the parameters and scope of the estimate}

\paragraph{Where the exponent arises.}
At exact density, the proof establishes
\[
 \delta_1(W,W_k)\leq A_{m,k}\Gamma+B_{m,k}\sigma,
 \qquad \sigma^2\leq\Gamma/c_m.
\]
The square root therefore enters when the squared degree deviation is
converted to $\sigma$. The spectral inequalities, thresholding,
extraction of the parts, and balancing are all bounded linearly in
$\Gamma$ and $\sigma$. In particular, the partition-recovery step
introduces no further loss in the exponent. Proposition~\ref{prop:sharpness}
shows that this square root is unavoidable for distance to the
balanced graphon $W_k$.

\paragraph{Growth with the cycle length.}
For fixed $k$, \eqref{eq:c-asymptotic} and the definitions imply
\[
 M_{m,k}=O_k\!\left(m(3/2)^{m/2}\right),
 \qquad
 \sqrt{L_{m,k}}=O_k\!\left(\sqrt m\,k^{m/2}\right).
\]
Thus
\begin{equation}\label{eq:C-growth}
 C_{m,k}=O_k\!\left(\sqrt m\,k^{m/2}\right).
\end{equation}
The exponent remains $1/2$ for every odd $m$.
For fixed $m,k$, the coefficient of $\sqrt{\delta}$ in the more
informative two-term bound
$L_{m,k}\delta+M_{m,k}\sqrt{\delta}$ is $M_{m,k}$.
The larger global coefficient $C_{m,k}$ makes a single-power bound
valid for every $\delta\geq0$.

There is an exponential lower bound for any such global coefficient.
Take the constant graphon $W\equiv p_k$, and write $q=1/k$, $p=p_k$.
Its cycle excess is
\[
 t(C_m,W)-G_m(p)=pq^{m-1}.
\]
For every relabelling of $W_k$, take $S=T$ equal to one of its parts.
On $S\times S$, $W_k^\phi=0$ whereas $W=p$, so
\[
 \dsq(W,W_k)\geq pq^2.
\]
A bound valid for all $\delta\geq0$ must therefore have
\begin{equation}\label{eq:C-lower}
 C(m,k)\geq\frac{pq^2}{\sqrt{pq^{m-1}}}
 =\sqrt p\,q^{(5-m)/2}
 =\sqrt{p_k}\,k^{(m-5)/2}.
\end{equation}
For fixed $k$, \eqref{eq:C-growth} and \eqref{eq:C-lower} have the
same exponential base $\sqrt{k}$; their ratio is at most of order
$\sqrt m$, with a multiplicative constant depending on $k$.
This lower bound concerns a bound valid for all $\delta$; it does not
assert the same restriction on a local coefficient when the permitted
range of $\delta$ may depend on $m$.

\paragraph{Degree control versus full structural control.}
The estimate \eqref{eq:degree-coercivity} controls the degree direction,
but not all departures from $W_k$ by itself.
For example, $W\equiv p_k$ has zero degree variance and
$\Phi_m(W)=0$, as follows directly from \eqref{eq:Phi-definition},
but its cut distance from $W_k$ is positive by the preceding test.
The complementary cycle--path defect in
\eqref{eq:gap-decomposition} supplies the additional information used
in Lemmas~\ref{lem:spectral} and \ref{lem:rounding}.

The balanced target is also essential when discussing possible linear
stability estimates. The graphons $W_\varepsilon$ in
Section~\ref{sec:sharpness} are already complete $k$-partite graphons.
They have distance zero from the family of all complete $k$-partite
graphons with arbitrary part measures, but distance of order
$\varepsilon$ from its balanced member. Our obstruction to a linear
rate concerns the latter distance. No linear estimate to the family
with unrestricted part measures is asserted here.

\paragraph{Nature of the proof.}
The qualitative compactness argument in Section~2 is not used in the
quantitative proof. The only smallness threshold in the constructive
partition argument is the explicit bound \eqref{eq:small-b}; above
that threshold the conclusion follows from the trivial bound
$\delta_1\leq1$. No regularity lemma or general graph removal lemma
is invoked.

\begin{thebibliography}{M}
\bibitem[M]{dense}
T.~Kim, J.~Baek, S.~Im, J.~Lee, and H.~Yang,
\emph{Goodman-type lower bounds for odd cycles I: The dense regime},
manuscript dated September 22, 2026.
\end{thebibliography}

\end{document}
